\Section{Свойства функций, непрерывных на отрезке — ограниченность, достижимость (точных) верхней и
нижней граней. Теорема о промежуточных значениях непрерывной функции. Теорема об обратной функции.}

\Subsection{Свойства функций, непрерывных на отрезке}

\begin{definition}
Функция \(f(x)\) называется \textit{непрерывной на отрезке \([a, \, b]\)}, если она
непрерывна в каждой внутренней точке отрезка, непрерывна справа в точке \(a\)
и непрерывна слева в точке \(b\).
\end{definition}

\begin{theorem}[Первая Теорема Вейерштрасса]
Непрерывная на отрезке \([a, \, b]\) функция \(f(x)\) ограничена на этом отрезке.

\begin{proof}
Пойдем от противного: пусть \(f(x)\) является неограниченной на отрезке \([a, \, b]\).
Тогда по определению неограниченности для любого натурального \(n\) найдется
такое число \(x_{n} \in [a, \, b]\), что \(|f(x_{n})| > n\). \\

Заметим, что из этого следует, что последовательность \(\{f(x_{n})\}\) является ББП.
Действительно, для любого \(\gamma > 0\) по аксиоме архимеда найдется натуральное \(N > \gamma\).
Но тогда для любых номеров \(n \geq N\) выполняется:
\[
|f(x_{n})| > n \geq N > \gamma
\]
\\
Последовательность \(\{x_{n}\}\) ограничена, поскольку все ее члены принадлежат
отрезку \([a, \, b]\). Тогда по теореме Больцано-Вейерштрасса существует
сходящаяся подпоследовательность \(\{x_{k_{n}}\}\) последовательности \(\{x_{n}\}\). \\

Пусть \(\lim\limits_{n \to \infty} x_{k_{n}} = c\). По Теореме \ref{2:theorem_lim_in_interval}
точка \(c\) принадлежит отрезку \([a, \, b]\). Тогда из непрерывности \(f(x)\) следует,
что \(\lim\limits_{n \to \infty} f(x_{k_{n}}) = f(c),\) то есть последовательность
\(\{f(x_{k_{n}})\}\) сходится. \\

Однако тогда мы получили два противоречащих факта:
\begin{itemize}
\item Подпоследовательность \(\{f(x_{k_{n}})\}\) последовательности \(\{f(x_{n})\}\)
сходится;
\item Подпоследовательность любой ББП является ББП (то есть не может сходиться).
\end{itemize}

Следовательно, наше предположение неверно, и теорема верна.
\end{proof}
\end{theorem}

\begin{theorem}[Вторая Теорема Вейерштрасса]
Непрерывная на отрезке \([a, \, b]\) функция \(f(x)\) достигает на этом отрезке
своих точных верхней и нижней граней.

\begin{proof}
Пойдем от противного: пусть \(\underset{x \in [a, \, b]}{\sup} f(x) = \alpha\) и
\(f(x) \neq \alpha\) при любых \(x \in [a, \, b]\).
Введем следующую функцию:
\[
g(x) = \frac{1}{\alpha - f(x)}.
\]
По свойствам непрерывных функций \(g(x)\) также будет непрерывной на отрезке \([a, \, b]\).
Тогда по первой теореме Вейерштрасса она же ограничена на этом отрезке некоторым \(\gamma > 0\). \\

Тогда справедливы неравенства (обратите внимание, что модуль слева не нужен):
\begin{gather}
    \frac{1}{\alpha - f(x)} \leq \gamma, \\
    f(x) \leq \alpha - \frac{1}{\gamma} < \alpha.
\end{gather}

Но тогда число \(\alpha\) не является точной верхней гранью~--- получили противоречие.
Следовательно, предположение неверно, и функция достигает своей точной верхней грани на отрезке. \\

Утверждение для точной нижней грани доказывается аналогично.
\end{proof}
\end{theorem}


\Subsection{Теорема о промежуточных значениях функции}

\begin{lemma} \label{6:lemma_sign_preservation}
Пусть функция \(f(x)\) определена в некоторой \(\Delta\)-окрестности точки \(x_{0}\).
Если функция \(f(x)\) непрерывна в точке \(x_{0}\) и \(f(x_{0}) \neq 0\),
то существует такая \(\delta\)-окрестность точки \(x_{0}\),
что функция \(f(x)\) сохраняет свой знак в этой окрестности.

\begin{proof}
Из непрерывности в точке \(x_{0}\) следует, что \(\lim\limits_{x \to x_{0}} f(x) = f(x_{0})\).
Тогда по определению Коши для любого \(\epsilon > 0\) найдется такое \(\delta > 0\),
что для любого \(x \in U_{\delta}(x_{0})\) выполняется:
\[
f(x_{0}) - \epsilon < f(x) < f(x_{0}) + \epsilon.
\]

Выберем \(\epsilon < |f(x_{0})|\). Тогда выражения \(f(x_{0}) - \epsilon\)
и \(f(x_{0}) + \epsilon\) принимают один знак, и тогда \(f(x)\) (по вышеуказанным неравенствам)
принимает тот же знак, что и \(f(x_{0})\)~--- что и требовалось доказать.
\end{proof}
\end{lemma}

\begin{note}
Если в доказанной лемме заменить обычную непрерывность на непрерывность слева (справа),
а обычную \(\delta\)-окрестность на левую (правую) \(\delta\)-окрестность,
то лемма останется верна.
\end{note}

\begin{theorem}[Теорема Больцано-Коши]
Если функция \(f(x)\) непрерывна на отрезке \([a, \, b]\)
и справедливо неравенство \(f(a) \cdot f(b) < 0\), то существует такая точка \(c \in [a, \, b]\),
что \(f(c) = 0\).

\begin{proof}
Неравенство \(f(a) \cdot f(b) < 0\) означает, что на концах отрезка функция принимает
значения разных знаков. Без ограничения общности будем считать, что \(f(a) < 0\) и \(f(b) > 0\). \\

Зададим множество \(X = {x \in [a, \, b]: f(x) < 0}\). Множество \(X\) заведомо непусто,
поскольку в него входит точка \(a\). Также \(X\) ограничено, поскольку \(X \subset [a, \, b]\).
Следовательно, по Теореме \ref{theorem_set_SupInf} существует \(\sup X = c\).
Докажем, что точка \(c\)~--- искомая. \\

По Теореме \ref{6:lemma_sign_preservation} найдутся такие \(\delta_{1} > 0\) и \(\delta_{2} > 0\),
что на интервале \([a, \, a + \delta_{1})\) функция принимает отрицательные значения,
а на полуинтервале \((b - \delta_{2}, \, b]\)~--- положительные.
Следовательно, имеем:
\begin{gather}
    c \in [a + \delta_{1}, \, b - \delta_{2}] \subset (a, \, b) \\
    c \in (a, \, b)
\end{gather}

Осталось доказать, что \(f(c) = 0\). \\
Предположим, что это не так и \(f(c) \neq 0\). По Теореме \ref{6:lemma_sign_preservation}
найдется \(\delta\)-окрестность точки \(c\), в которой функция \(f(x)\) сохраняет свой знак.
По определению точной верхней грани найдется такая точка \(x' \in X \cap (c - \delta, \, c)\),
что \(f(x') < 0\). С другой стороны, для всех точек из интервала \((c, \, c + \delta)\)
выполняется неравенство \(f(x) \geq 0\). Но тогда знак функции в \(\delta\)-окрестности
не сохраняется~--- получили противоречие, следовательно, наше предположение
неверно и \(f(c) = 0\).
\end{proof}
\end{theorem}

\begin{theorem}[Теорема о промежуточных значениях функции]
Пусть функция \(f(x)\) непрерывна на отрезке \([a, \, b]\) и \(f(a) \neq f(b)\).
Тогда для любого числа \(\gamma\), лежащего между значениями \(f(a)\) и \(f(b)\), найдется
такая точка \(c \in (a, \, b)\), что \(f(c) = \gamma\).

\begin{proof}
Для определенности будем считать, что \(f(a) > f(b)\), и тогда \(f(a) > \gamma > f(b)\).
Рассмотрим функцию \(g(x) = f(x) - \gamma\). Она непрерывна на отрезке \([a, \, b]\),
и притом \(g(a) = f(a) - \gamma > 0\) и \(g(b) = f(b) - \gamma < 0\).
Тогда по Теореме Больцано-Коши найдется такая точка \(c \in (a, \, b)\), что \(g(c) = 0\).
Следовательно, \(f(c) = g(c) + \gamma = \gamma\), и точка \(c\)~--- искомая.
\end{proof}
\end{theorem}


\Subsection{Теорема об обратной функции}

\begin{definition}
Пусть на множествах \(X \subset \R\) и \(Y \subset \R\) задана функция \(f:X \to Y\).
Предположим, что функция \(y = f(x)\) обладает следующим свойством:
каждому \(y_{0} \in Y\) соответствует \textit{единственное} число \(x_{0} \in X\) такое,
что \(f(x_{0}) = y_{0}.\) \\

Тогда на множестве \(Y\) можно определить \textit{обратную функцию} \(x = f^{-1} (y)\),
ставя в соответствие каждому числу \(y \in Y\) такое число \(x \in X\),
что \(y = f(x)\). \\

Функция \(y = f(x)\) называется \textit{обратимой}. \\
Функция \(y = f(x)\) является обратной к функции \(x = f^{-1} (y)\), поэтому
\(y = f(x)\) и \(x = f^{-1} (y)\) называют \textit{взаимно обратными функциями}.
\end{definition}

\begin{lemma} \label{6:lemma_inverse_function}
Если функция строго монотонная на множестве \(X\), а \(Y\)~--- множество ее значений,
то на множестве \(Y\) определена обратная функция \(x = f^{-1} (y)\). \\
Более того, если \(y = f(x)\) возрастающая (убывающая) на множестве \(X\),
то функция \(x = f^{-1} (y)\) также является возрастающей (убывающей) на множестве \(Y\).

\begin{proof}
Поскольку функция \(y = f(x)\) является строго монотонной на \(X\),
то она удовлетворяет условию взаимно однозначного соответствия аргументов и значений функции.
Действительно: для любых \(x_{1} \neq x_{2}\) из множества \(X\) значения функции
\(f(x_{1})\) и \(f(x_{2})\) не равны, ведь иначе нарушалась бы строгая монотонность.
Следовательно, по вышеописанному определению существует функция \(x = f^{-1} (y),\)
являющаяся обратной к данной. \\

Вторую часть леммы докажем для случая возрастающей функции (случай убывающей аналогичен).
Пусть функция \(y = f(x)\) является возрастающей на \(X\). Для возрастания \(x = f^{-1} (y)\)
нам нужно доказать, что для любых \(y_{1} > y_{2}\) из множества \(Y\) выполняется:
\[
f^{-1} (y_{1}) = x_{1} > x_{2} = f^{-1} (y_{2}).
\]
Предположим противное: пусть \(x_{1} \leq x_{2}\). Тогда из монотонности исходной функции
следует, что \(y_{1} \leq y_{2}\). Однако по условию \(y_{1} > y_{2}\)~--- противоречие.
Следовательно, предположение неверно, и тогда лемма верна.
\end{proof}
\end{lemma}

\begin{theorem}
Если функция \(y = f(x)\) возрастает (убывает) и непрерывна на отрезке \([a, \, b]\),
кроме того, \(f(a) = \alpha\) и \(f(b) = \beta\), то на множестве \([\alpha, \, \beta]\)
(на множестве \([\beta, \, \alpha]\)) определена и непрерывна обратная функция \(x = f^{-1} (y)\),
которая возрастает (убывает) на этом множестве.

\begin{proof}
В Лемме \ref{6:lemma_inverse_function} мы уже доказали существование и строгую монотонность
обратной функции \(x = f^{-1} (y)\)~--- следовательно, осталось доказать ее непрерывность. \\

Проведем доказательство для возрастающей функции (для убывающей~--- аналогично).
Пусть \(x_{0} \in (a, \, b)\) и \(y_{0} = f(x_{0})\).
Заметим, что если мы докажем равенство:
\[
f^{-1} (y_{0} - 0) = f^{-1} (y_{0}) = f^{-1} (y_{0} + 0),
\]
то отсюда (см. свойства односторонних пределов) будет следовать, что
\[
\exists \, \lim_{ y \to y_{0} } f^{-1} (y) = f^{-1} (y_{0}),
\]
то есть по определению функция будет непрерывной в точке. \\

Докажем левое равенство.
Предположим противное, то есть что левое равенство не выполнено.
Тогда по Теореме о существовании односторонних пределов у монотонных функций получаем:
\[
f^{-1} (y_{0} - 0) < f^{-1} (y_{0}).
\]
Заметим, что для любой точки \(y \in [\alpha, \, y_{0})\) выполнено:
\[
a \leq f^{-1} (y) \leq f^{-1} (y_{0} - 0),
\]
а для любой точки \(y \in [y_{0}, \, \beta]\) справедливы неравенства:
\[
f^{-1} (y_{0}) \leq f^{-1} (y) \leq b.
\]

Но тогда мы получаем, что интервал \((f^{-1} (y_{0} - 0), \, f^{-1} (y_{0}))\)
не принадлежит множеству значений функции \(x = f^{-1} (y)\). Однако это противоречит тому,
что для \textit{непрерывной} функции \(y = f(x)\) образом отрезка \([a, \, b]\)
является отрезок \([\alpha, \, \beta]\). Следовательно, предположение неверно,
и \(f^{-1} (y_{0} - 0) = f^{-1} (y_{0}).\) \\

Аналогично доказывается, что \(f^{-1} (y_{0})  = f^{-1} (y_{0} + 0)\).
Следовательно, функция \(x = f^{-1} (y)\) является непрерывной на интервале \((\alpha, \, \beta)\).
Непрерывность справа и слева в соответствующих концах отрезка доказывается аналогично.
\end{proof}
\end{theorem}

\begin{center}
\begin{tikzpicture}[x=0.8cm,y=0.8cm,>=stealth]

% Верхняя числовая прямая
\draw[->,black] (0.5,6.5) -- (15.5,6.5);
\draw[red,line width=2.5pt] (2,6.5) -- (14,6.5);
\draw[black] (2,6.35) -- (2,6.65);
\draw[black] (14,6.35) -- (14,6.65);
\node[below] at (2,6.2) {$a$};
\node[below] at (14,6.2) {$b$};
\node[below] at (15.5,6.4) {$x$};

% Стрелка к средней прямой
\draw[->,black] (8,5.775) -- (8,3.725);
\node[right] at (8.2,4.75) {$y=f(x)$};

% Средняя числовая прямая
\draw[->,black] (0.5,3) -- (15.5,3);
\draw[red,line width=2.5pt] (3.5,3) -- (13,3);
\draw[black] (3.5,2.85) -- (3.5,3.15);
\draw[black] (10,2.85) -- (10,3.15);
\draw[black] (13,2.85) -- (13,3.15);
\node[below] at (3.5,2.7) {$\alpha$};
\node[below] at (10,2.7) {$y_0$};
\node[below] at (13,2.7) {$\beta$};
\node[below] at (15.5,2.9) {$y$};

% Стрелка к нижней прямой
\draw[->,black] (8,2.075) -- (8,0.225);
\node[right] at (8.2,1.25) {$x=f^{-1}(y)$};

% Нижняя числовая прямая
\draw[->,black] (0.5,-0.5) -- (15.5,-0.5);

% Тёмно-серая штриховка
\path[pattern=north east lines,pattern color=black!70]
    (6.4,-0.66) rectangle (10.5,-0.34);

% Красные отрезки
\draw[red,line width=2.5pt] (2,-0.5) -- (6.4,-0.5);
\draw[red,line width=2.5pt] (10.5,-0.5) -- (14,-0.5);

% Отмеченные точки нижней прямой
\foreach \x in {2,6.4,10.5,14} {
    \draw[black] (\x,-0.65) -- (\x,-0.35);
}

\node[below] at (2,-0.7) {$a$};
\node[below] at (6.4,-0.7) {$f^{-1}(y_0-0)$};
\node[below] at (10.5,-0.7) {$f^{-1}(y_0)$};
\node[below] at (14,-0.7) {$b$};
\node[below] at (15.5,-0.5) {$x$};

\end{tikzpicture}
\\
\textbf{Противоречие с <<выпавшим>> интервалом}
\end{center}